\documentclass[slidetop,11pt]{beamer}
%
% Ces deux lignes {\`a} d{\'e}commenter pour sortir 
% le texte en classe article
% \documentclass[class=article,11pt,a4paper]{beamer}
% \usepackage{beamerbasearticle}

% Packages pour les fran\c{c}ais
%
\usepackage[T1]{fontenc} 
\usepackage[latin1]{inputenc}
\usepackage[frenchb]{babel}
% pour un pdf lisible {\`a} l'{\'e}cran si on ne dispose pas 
% des fontes cmsuper ou lmodern
%\usepackage{lmodern}
\usepackage{aeguill}

% Pour afficher le pdf en plein ecran
% (comment{\'e} pour imprimer les transparents et pour les tests)
%\hypersetup{pdfpagemode=FullScreen}

% Red{\'e}finit la couleur de fond pour imprimer sur transparents
%\xdefinecolor{fondtexte}{rgb}{1,1,1}     % blanc

% commande differente pour les couleurs nomm{\'e}es - de base
%\colorlet{coultexte}{black} 

% -------------- Fioritures de style -------------
% Fait afficher l'ensemble du frame 
% en peu lisible (gris clair) d{\`e}s l'ouverture
\beamertemplatetransparentcovered

% Supprimer les icones de navigation (pour les transparents)
\setbeamertemplate{navigation symbols}{}

% Mettre les icones de navigation en mode vertical (pour projection)
% \setbeamertemplate{navigation symbols}[vertical]

\usetheme{Warsaw}

%------------ fin style beamer -------------------

% Faire appara{\^i}tre un sommaire avant chaque section
% \AtBeginSection[]{
%   \begin{frame}
%   \frametitle{Plan}
%   \medskip
%   %%% affiche en d{\'e}but de chaque section, les noms de sections et
%   %%% noms de sous-sections de la section en cours.
%   \small \tableofcontents[currentsection, hideothersubsections]
%   \end{frame} 
% }

% ----------- Contenu de la page de titre --------
\title{Recruter en informatique}
\subtitle{Les mauvaises id{\'e}es... et les bonnes !} 
\author{Gaby Wald}
\institute{\emph{to be defined}}
%% \date{\today} \date{19 mai 2020} %% \date{\today} %% \date{03 Janvier 2018}
\logo{\includegraphics[height=0.5cm]{img/logo_glider.png}}

% ----------- Nom des diff{\'e}rentes parties --------

%% Mis au fur et {\`a} mesure...

% ------------------------------------------------
% -------------   D{\'e}but document   ---------------
% ------------------------------------------------
\begin{document}
%--------- {\'e}criture de la page de titre ----------
% avec la commande frame simplifi{\'e}e
\frame[plain]{\titlepage } 
%

%------------------ Sommaire ---------------

\begin{frame}
	\frametitle{Sommaire}
	\small \tableofcontents[hideallsubsections]
\end{frame} 

\section{Recrutement en informatique, une introduction}
\begin{frame}
	\frametitle{Recrutement en informatique, une introduction}
	\tableofcontents[sections=1,currentsection,subsectionstyle=show/shaded/hide] %% sectionstyle=hide/hide,
\end{frame}

\subsection{Introduction}
\begin{frame}
	\frametitle{Introduction}
	\begin{itemize}
		\item L'objectif ici est seulement de recenser une liste de pratiques ; 
		\item Enlever les clich{\'e}s RH sur les infomaticiens, d{\'e}veloppeurs, ing{\'e}nieurs... ; 
		\item Faire sauter les soucis principaux du recrutement (nombre de candidats, turn-over, refus de contact...) ;
		\item Faire {\'e}voluer les mentalit{\'e}s (divas des deux bords, sexisme...) ; 
		\item ...  
	\end{itemize}
\end{frame}

\subsection{Mauvaises id{\'e}es, pr{\'e}jug{\'e}s...}
\begin{frame}
	\frametitle{Mauvaises id{\'e}es, pr{\'e}jug{\'e}s...}
	\begin{itemize}
		\item Recrutement uniquement sur mots-clefs et ignorer le contenu ; 
		\item Que des femmes pour recruter des hommes ("LinkedIn est un Tinder invers{\'e}. ") ; 
		\item Recrutement sur "gadget" : comp{\'e}tences, ann{\'e}es d'exp{\'e}riences... et le sens ?
		\item Recherche de "M{\^a}les Alphas C{\'e}libataires Sans Enfants" ; 
		\item R{\'e}f{\'e}rences ? Quelles r{\'e}f{\'e}rences ? Pourquoi ?
		\item Pr{\'e}sentation g{\'e}n{\'e}rique, "recrutement sur profil", "grille de salaire"... ; 
		\item "Le digital : {\`a} manipuler avec doigt{\'e} !" : num{\'e}rique !
		\item Et la culture associ{\'e}e... 
	\end{itemize}
\end{frame}

\subsection{Bonnes id{\'e}es (?)}
\begin{frame}
	\frametitle{Bonnes id{\'e}es (?)}
	\begin{itemize}
		\item Instruisez-vous, renseignez-vous, oubliez vos clich{\'e}s...
		\item Communiquer directement les enjeux (missions, r{\'e}mun{\'e}ration, conditions de travail...) ;
		\item Environnement bienveillant, tol{\'e}rant, ouvert, inclusif, divers...
		\item Distinction Vie Pro et Vie Perso...
		\item Reconnaissance et assurer l'avenir : {\^e}tre correct avec tout le monde, y compris les contractuels / prestataires / stagiaires...
		\item Une {\'e}quipe, cela se construit, et il n'y a pas d'argent magique, ni de management magique...
		%% \item ...
	\end{itemize}
	
	\emph{Note : une obligation l{\'e}gale sociale n'est pas un avantage (mutuelle, transports...). } 
\end{frame}

%%%%% %%%%% %%%%% %%%%% %%%%% 

\section{Recrutement en informatique : les mauvaises pratiques}
\begin{frame}
	\frametitle{Recrutement en informatique : les mauvaises pratiques}
	\tableofcontents[sections=2,currentsection,subsectionstyle=show/shaded/hide] %% sectionstyle=hide/hide,
\end{frame}

\subsection{Pr{\'e}jug{\'e}s "techniques"}
\begin{frame}
	\frametitle{Pr{\'e}jug{\'e}s "techniques"}
	\begin{itemize}
		\item Aspects techniques peu attirants... ; 
		\item Milieu Geek / Nerd / "Alternatif"... ; 
		\item Terminologie magique ("digital / num{\'e}rique", acronymes, "Java / JavaScript"...) ; 
		\item Informatique magique ("c'est simple", parfois non) ;  
		\item[] 
		\item Comme tout, \c{c}a s'apprend ! Au moins une culture minimale sur les acronymes... 
	\end{itemize}
\end{frame}

\subsection{Soucis de recrutement}
\begin{frame}
	\frametitle{Soucis de recrutement}
	\begin{itemize}
		\item Turn-Over, concurrence des recruteurs (ESN / SS2I), harc{\`e}lement de contact... 
		\item Recrutement uniquement sur mots-clefs et ignorer le contenu ; 
		\item Que des femmes pour recruter des hommes ("LinkedIn est un Tinder invers{\'e}. ") ; 
		\item Recrutement sur "gadget" : comp{\'e}tences, ann{\'e}es d'exp{\'e}riences... et le sens ?
		\item Recherche de "M{\^a}les Alphas C{\'e}libataires Sans Enfants" ; ...
	\end{itemize}
\end{frame}

\subsection{Lev{\'e}es de drapeaux rouges !}
\begin{frame}
	\frametitle{Lev{\'e}es de drapeaux rouges !}
	\begin{itemize}
		\item Soci{\'e}t{\'e} de Service en Informatique ("A*" / externalisations) ; 
		\item Pr{\'e}sentation comme avantages ce qui rel{\`e}ve du l{\'e}gal (cong{\'e}s pay{\'e}s, tickets restaurants, remboursement des transports...) ; 
		\item Ambiance(s) toxique(s) : management inadapt{\'e}, "mascus et alphas", prestataires / consultants...
		\item Recrutement sur profil (bloquer la concurrence ?) ; 
		\item Recrutement en comp{\'e}tition (fa\c{c}on consultant commercial) ; 
		\item Conditions : OpenSpace avec plus de 10 personnes, T{\'e}l{\'e}Travail limit{\'e}...
		\item ... 
	\end{itemize}
\end{frame}

%%%%% %%%%% %%%%% %%%%% %%%%% 

\section{Recrutement en informatique : pour faire mieux}
\begin{frame}
	\frametitle{Recrutement en informatique : pour faire mieux}
	\tableofcontents[sections=3,currentsection,subsectionstyle=show/shaded/hide] %% sectionstyle=hide/hide,
\end{frame}

\subsection{Bonnes id{\'e}es (1)}
\begin{frame}
	\frametitle{Bonnes id{\'e}es (1)}
	\begin{itemize}
		\item Instruisez-vous, renseignez-vous, oubliez vos clich{\'e}s...
		\item Communiquer directement les enjeux (missions, r{\'e}mun{\'e}ration, conditions de travail...) ;
		\item Environnement bienveillant, tol{\'e}rant, ouvert, inclusif, divers...
		\item Distinction Vie Pro et Vie Perso...
		\item Reconnaissance et assurer l'avenir : {\^e}tre correct avec tout le monde, y compris les contractuels / prestataires / stagiaires...
		\item Une {\'e}quipe, cela se construit, et il n'y a pas d'argent magique, ni de management magique...
		\item Aller chercher les candidats (et s'int{\'e}resser {\`a} leur CV)...
	\end{itemize} 
\end{frame}

\subsection{Bonnes id{\'e}es (2)}
\begin{frame}
	\frametitle{Bonnes id{\'e}es (2)}
	\begin{itemize}
		\item ... 
	\end{itemize} 
\end{frame}

\subsection{Bonnes id{\'e}es (3)}
\begin{frame}
	\frametitle{Bonnes id{\'e}es (3)}
	\begin{itemize}
		\item ... 
	\end{itemize} 
\end{frame}

%% \def\sectionPartBibliographie{Bibliographie / Mediagraphie}
%% \section{\sectionPartBibliographie}
%% \begin{frame}
%% 	\frametitle{\sectionPartBibliographie}
%% 	{\footnotesize 
%% 		\nocite{*}
%% 		%toutes references biblio : 6 lettres + 2 chiffres
%% 		\bibliography{presentationRecrutementInformatique}
%% 		% \bibliographystyle{frplain} % plain or frplain
%% 		\bibliographystyle{plain}
%% 	}
%% \end{frame}

\end{document}
